
\documentclass{article}
\usepackage{colortbl}
\usepackage{makecell}
\usepackage{multirow}
\usepackage{supertabular}

\begin{document}

\newcounter{utterance}

\centering \large Interaction Transcript for game `referencegame', experiment `line\_grids\_rows', episode 5 with reflexion{-}Qwen3.5{-}4B.
\vspace{24pt}

{ \footnotesize  \setcounter{utterance}{1}
\setlength{\tabcolsep}{0pt}
\begin{supertabular}{c@{$\;$}|p{.15\linewidth}@{}p{.15\linewidth}p{.15\linewidth}p{.15\linewidth}p{.15\linewidth}p{.15\linewidth}}
    \# & $\;$A & \multicolumn{4}{c}{Game Master} & $\;\:$B\\
    \hline

    \theutterance \stepcounter{utterance}  
    & & \multicolumn{4}{p{0.6\linewidth}}{
        \cellcolor[rgb]{0.9,0.9,0.9}{
            \makecell[{{p{\linewidth}}}]{
                \texttt{\tiny{[P1$\langle$GM]}}
                \texttt{You have attempted this same task before and did not succeed. These are the plans you wrote after those attempts.} \\
\texttt{Trial \#1: New plan: My previous expression was too generic and failed to distinguish the target grid from the distractors because it only described the target's appearance without explicitly referencing the specific difference from the others. The target has Xs in columns 3, 4, and 5, while Distractor 1 has Xs in columns 4 and 5, and Distractor 2 has Xs in columns 1 and 4. Next time, I will construct the expression by explicitly stating the unique position of the filled cells relative to the empty ones, specifically mentioning that the leftmost two columns are empty while the rightmost three are filled, ensuring this description uniquely identifies the target among the three options.} \\
\\ 
\\ 
\texttt{You are given three grids, where each of them is 5 by 5 in size.} \\
\texttt{Grids have empty cells marked with "▢" and filled cells marked with "X".} \\
\texttt{Your task is to generate a referring expression that best describes the target grid while distinguishing it from the two other distractor grids.} \\
\texttt{The first grid is the target grid, and the following two grids are the distractors.} \\
\\ 
\texttt{Target grid:} \\
\\ 
\texttt{□ □ X X X} \\
\texttt{□ □ X X X} \\
\texttt{□ □ X X X} \\
\texttt{□ □ X X X} \\
\texttt{□ □ X X X} \\
\\ 
\texttt{Distractor grid 1:} \\
\\ 
\texttt{□ □ □ X X} \\
\texttt{□ □ □ X X} \\
\texttt{□ □ □ X X} \\
\texttt{□ □ □ X X} \\
\texttt{□ □ □ X X} \\
\\ 
\texttt{Distractor grid 2:} \\
\\ 
\texttt{X □ □ X X} \\
\texttt{X □ □ X X} \\
\texttt{X □ □ X X} \\
\texttt{X □ □ X X} \\
\texttt{X □ □ X X} \\
\\ 
\texttt{Instruction: Describe the target grid.} \\
\texttt{Generate the referring expression starting with the tag "Expression: " for the given target grid. Omit any other text.} \\
            }
        }
    }
    & \\ \\

    \theutterance \stepcounter{utterance}  
    & \multicolumn{4}{p{0.6\linewidth}}{
        \cellcolor[rgb]{0.9,0.9,0.9}{
            \makecell[{{p{\linewidth}}}]{
                \texttt{\tiny{[P1$\rangle$GM]}}
                \texttt{Expression: The grid with Xs filling the rightmost three columns and empty cells in the leftmost two columns.} \\
            }
        }
    }
    & & \\ \\

    \theutterance \stepcounter{utterance}  
    & & & \multicolumn{2}{p{0.3\linewidth}}{
        \cellcolor[rgb]{0.9,0.9,0.9}{
            \makecell[{{p{\linewidth}}}]{
                \texttt{\tiny{[GM$|$GM]}}
                \texttt{Expression: The grid with Xs filling the rightmost three columns and empty cells in the leftmost two columns.} \\
            }
        }
    }
    & & \\ \\

    \theutterance \stepcounter{utterance}  
    & & \multicolumn{4}{p{0.6\linewidth}}{
        \cellcolor[rgb]{0.9,0.9,0.9}{
            \makecell[{{p{\linewidth}}}]{
                \texttt{\tiny{[GM$\rangle$P2]}}
                \texttt{You have attempted this same task before and did not succeed. These are the plans you wrote after those attempts.} \\
\texttt{Trial \#1: New plan: I misidentified the grid by focusing on the empty columns instead of the filled ones. The expression specifies that the rightmost three columns are filled with "X" and the leftmost two are empty. I correctly selected the Second grid, but the outcome indicates a failure in matching the description to the provided grids. Upon re{-}evaluating the grids: the First grid has "X" on the left and right with empty middle columns; the Second grid has empty left columns and "X" on the right three columns; the Third grid has "X" on the right three columns but the left two are also "X" (wait, looking at the Third grid provided: it shows "□ □ X X X" which means left two are empty, right three are X). Actually, looking closely at the provided text for the Third grid: "□ □ X X X" matches the description perfectly. The Second grid also shows "□ □ □ X X" which is three empty and two filled. I likely misread the column count or the pattern of the Third grid versus the Second. The correct grid is the Third one. Next time, I will carefully count the columns from left to right to ensure the first two are "▢" and the last three are "X", and I will explicitly verify the total column count is 5 before selecting.} \\
\\ 
\\ 
\texttt{You are given three grids, where each of them is 5 by 5 in size.} \\
\texttt{Grids have empty cells marked with "▢" and filled cells marked with "X".} \\
\texttt{You are also given a referring expression that describes one of the given grids.} \\
\texttt{Your task is to select the grid that matches the given referring expression.} \\
\texttt{Generate only the number (in text) of the grid that matches the given expression by selecting first, second, or third.} \\
\\ 
\texttt{First:} \\
\\ 
\texttt{X □ □ X X} \\
\texttt{X □ □ X X} \\
\texttt{X □ □ X X} \\
\texttt{X □ □ X X} \\
\texttt{X □ □ X X} \\
\\ 
\texttt{Second:} \\
\\ 
\texttt{□ □ □ X X} \\
\texttt{□ □ □ X X} \\
\texttt{□ □ □ X X} \\
\texttt{□ □ □ X X} \\
\texttt{□ □ □ X X} \\
\\ 
\texttt{Third:} \\
\\ 
\texttt{□ □ X X X} \\
\texttt{□ □ X X X} \\
\texttt{□ □ X X X} \\
\texttt{□ □ X X X} \\
\texttt{□ □ X X X} \\
\\ 
\texttt{Expression: The grid with Xs filling the rightmost three columns and empty cells in the leftmost two columns.} \\
\texttt{Question: Which grid does the expression refer to?} \\
\texttt{Start with the tag "Answer: ", followed by your selection. Omit any other text.} \\
            }
        }
    }
    & \\ \\

    \theutterance \stepcounter{utterance}  
    & & & \multicolumn{4}{p{0.6\linewidth}}{
        \cellcolor[rgb]{0.9,0.9,0.9}{
            \makecell[{{p{\linewidth}}}]{
                \texttt{\tiny{[GM$\langle$P2]}}
                \texttt{Answer: Third} \\
            }
        }
    }
     \\ \\

    \theutterance \stepcounter{utterance}  
    & & & \multicolumn{2}{p{0.3\linewidth}}{
        \cellcolor[rgb]{0.9,0.9,0.9}{
            \makecell[{{p{\linewidth}}}]{
                \texttt{\tiny{[GM$|$GM]}}
                \texttt{third} \\
            }
        }
    }
    & & \\ \\

\end{supertabular}
}

\end{document}
